\batchmode
\documentstyle[11pt]{article}
\addtolength{\textheight}{1.5in}
\addtolength{\textwidth}{1.5in}
\addtolength{\oddsidemargin}{-.5in}
\addtolength{\evensidemargin}{-.5in}
\begin{document}
\begin{center}
\begin{LARGE}
Categorical Programming Language \\[1em]
\end{LARGE}
September 1985 \\[1em]
Tatsuya Hagino \\[1em]
\begin{small}
Department of Computer Science \\
University of Edinburgh \\
James Clerk Maxwell Building \\
The King's Buildings \\
Mayfield Road \\
Edinburgh EH9 3JZ
\end{small}
\end{center}

\section{Introduction}

This note briefly explains what {\it Categorical Programming Language}
(`CPL' for short) is and its system which is implemented on VAX UNIX 4.2.
One of the main features of CPL is its declaration mechanism for data types.

\begin{enumerate}
\item CPL uses categorical characterizations of data types.  For example,
product $A\times B$ is characterized by giving two units, $\pi_{1}{\rm
:}A\times B\longrightarrow A$ and $\pi_{2}{\rm :}A\times B\longrightarrow
B$, having the universal property that for any two morphisms, $f{\rm
:}C\longrightarrow A$ and $g{\rm :}C\longrightarrow B$, there exists a
unique morphism $h {\rm :}C\longrightarrow A\times B$ such that $\pi_{1}\circ
h = f$ and $\pi_{2}\circ h = g$.

\item CPL has no primitive data types.  Ordinary programming languages
(e.g.\ Pascal, Lisp, ML) have primitive data types: natural number, list,
record (i.e.\ product), disjoint record (i.e.\ coproduct) and so on, but CPL
can define these data types.  Thus, CPL is analogous to algebraic
specification methods which can specify algebras like natural number and
list.  However they cannot specify higher order data types (e.g.\
exponentiation) easily, nor can they specify product without using equations
(connecting projection and pairing operations).

\item CPL can not only define a product without equations, but also can
define a function space which is categorically characterized as
exponentiation.  Furthermore, although its use is still questionable, right
adjoint of exponentiation can be defined as well.

\item CPL is symmetric in the sense that it allows one to define data types by
both initial and terminal data constraints.  This is not the case with
algebraic specification methods.

\item As a result of this symmetry, CPL can define final co-algebras, which
includes objects like infinite lists.

\item CPL defines functors.  Algebraic specification methods usually
introduce parametrization later, but in CPL functorial behaviour of
parametrized data types plays an important roll.
\end{enumerate}

Algebraic specification methods specify algebras, whereas CPL specifies
categories.  A CPL specification consists of a series of object
declarations.  Each object declaration specifies a functor, some associated
natural transformations and some universal property.  A model of a CPL
specification is a category with the specified functors and natural
transformations satisfying the specified universal properties.  For example,
if a specification consists of one object declaration which introduces a
binary functor `prod' with two natural transformations (i.e.\ its
projections) and the universal property of which is that {\it `prod' is the
right adjoint functor of the diagonal functor}, a model of the specification
is a category ${\cal C}$ with product and `prod' should be interpreted as
the product functor ${\cal C} \times {\cal C} \longrightarrow {\cal C}$.

By defining suitable morphisms between models, models of a specification
form a category.  If a specification is enriched by object declarations, the
category of models of the enriched specification is a sub-category of the
category of models of the original specification.  When CPL starts, an empty
specification is given, so that the category of models is the category of
categories, and by declaring objects the category of models becomes smaller
and smaller, e.g.\ the category of categories, the category of
product categories, the category of cartesian closed categories, $\ldots$.

A simplified object declaration in CPL looks like:
\begin{quote}
right object $O(a)$ with $\psi$ is \\
\makebox[2em]{} $u$ : $F(a,O)$ $\longrightarrow$ $G(a,O)$ \\
end object \\
\end{quote}
where $F$ and $G$ are binary functors which have already been specified.
The specification is enriched with a unary functor $O$ and a natural
transformation $u$ from $F(a,O(a))$ to $G(a,O(a))$ (we call $u$ unit
morphism).  The universal property is: for any morphism $f{\rm
:}F(a,b)\longrightarrow G(a,b)$, there exists a unique morphism $h {\rm :} b
 \longrightarrow O(a)$ such that the following diagram commutes:
\begin{displaymath}
\begin{array}{rcccl}
& & f & & \\
& F(a,b) & \longrightarrow & G(a,b) & \\
F(a,h) & {\Big\downarrow} & & {\Big\downarrow} & G(a,h) \\
& F(a,O(a)) & \longrightarrow & G(a,O(a)) & \\
& & u & & \\
\end{array}
\end{displaymath}
For the moment we have assumed that $F$ and $G$ are covariant in their second
argument.  $\psi$ in the declaration is called the `factorizer' and is used
to denote the unique arrow, i.e.\ $h$ is denoted as $\psi (f)$.

The universal property is equivalent to the following equations:
\begin{displaymath}
\begin{array}{l}
\psi (u) = id_{O(a)} \\
u\circ F(a,\psi (f)) = G(a,\psi (f))\circ f \\
\psi (f)\circ g = \psi (h) \\
\makebox[5em]{} {\rm where} \; f\circ F(a,g) = G(a,g)\circ h \\
O(f) = \psi (G(f,O(a))\circ u\circ F(f,O(a)) \\
\end{array}
\end{displaymath}
The functor $O$ may not necessary be covariant in its argument.  It depends
on the variance of $F$ and $G$.

In the declaration above, we used keyword `right', which indicates that $O$
has characteristics of right adjoint or the terminal object.  We can replace
`right' with `left' getting a functor which has characteristics of left
adjoint or the initial object.
\begin{quote}
left object $\hat{O}(a)$ with $\hat{\psi}$ is \\
\makebox[2em]{} $\hat{u}$ : $\hat{F}(a,\hat{O})$ $\longrightarrow$
	$\hat{G}(a,\hat{O})$ \\
end object \\
\end{quote}
For any morphism $f{\rm :}\hat{F}(a,b)\longrightarrow \hat{G}(a,b)$, there
exists a unique morphism denoted by $\hat{\psi}(f)$ from $\hat{O}(a)$ to $b$
such that the following diagram commutes:
\begin{displaymath}
\begin{array}{rcccl}
& & \hat{u} & & \\
& \hat{F}(a,\hat{O}(a)) & \longrightarrow & \hat{G}(a,\hat{O}(a)) & \\
\hat{F}(a,\hat{\psi}(f)) & {\Big\downarrow} & & {\Big\downarrow} &
	\hat{G}(a,\hat{\psi}(f)) \\
& \hat{F}(a,b) & \longrightarrow & \hat{G}(a,b) & \\
& & f & & \\
\end{array}
\end{displaymath}
The equations equivalent to the universal property are:
\begin{displaymath}
\begin{array}{l}
\hat{\psi}(\hat{u}) = id_{\hat{O}(a)} \\
\hat{G}(a,\hat{\psi}(f))\circ \hat{u} = f\circ \hat{F}(a,\hat{\psi}(f)) \\ 
g\circ \hat{\psi}(f) = \hat{\psi}(h) \\
\makebox[5em]{} {\rm where} \; \hat{G}(a,g)\circ f = h\circ \hat{F}(a,g) \\
\hat{O}(f) = \hat{\psi}(\hat{G}(f,\hat{O}(a))\circ \hat{u}\circ
	\hat{F}(f,\hat{O}(a)) \\
\end{array}
\end{displaymath}

In general an object declaration in CPL can define a functor of many
variables (including nullary functors) and can also have any number of unit
morphisms.


\section{Examples of Object Declarations}

Since there are no primitive data types pre-declared in CPL, the first thing
we normally have to do is to define the terminal object.  It can be done by
declaring `\verb"1"' as follows.
\begin{quote}
\begin{verbatim}
right object 1 with !
end object
\end{verbatim}
\end{quote}
It does not have any unit morphisms.  By the characterization in the
previous section, there exists a unique morphism denoted by `\verb"!"' from
any object $a$ to `\verb"1"'.  Therefore `\verb"1"' really specifies the
terminal object.

Next, we define product by:
\begin{quote}
\begin{verbatim}
right object prod(a,b) with pair is
  pi1: prod -> a
  pi2: prod -> b
end object
\end{verbatim}
\end{quote}
Note that in this case $F$ and $G$ are projection functors.  There are two
unit morphisms, \verb"pi1:prod(a,b)->a" and \verb"pi2:prod(a,b)->b".  For
any morphisms \verb"f:c->a" and \verb"g:c->b" factorizer `\verb"pair"' gives
a unique morphism \verb"pair(f,g):c->prod(a,b)" such that the following
equations hold.
\begin{quote}
\begin{verbatim}
pair(pi1,pi2) = id
pi1.pair(f,g) = f
pi2.pair(f,g) = g
pair(f,g).h = pair(f.h,g.h)
prod(f,g) = pair(f.pi1,g.pi2)
\end{verbatim}
\end{quote}
Note that `\verb"id"' is the identity morphism and \verb"f.g" denotes the
composition of morphisms \verb"f" and \verb"g" (for \verb"f:b->c" and
 \verb"g:a->b", \verb"f.g:a->c".  Alternatively it can be written as
 \verb"g|f").

In order to do certain calculations, categories often need to be cartesian
closed.  We define exponentiation in CPL as follows:
\begin{quote}
\begin{verbatim}
right object exp(a,b) with curry
  ev: prod(exp,a) -> b
end object
\end{verbatim}
\end{quote}
For any morphism \verb"f:prod(c,a)->b", \verb"curry(f)" is from \verb"c" to
 \verb"exp(a,b)", and the following equations hold:
\begin{quote}
\begin{verbatim}
curry(ev) = id
ev.prod(curry(f),id) = f
curry(f).g = curry(f.prod(id,g))
exp(f,g) = curry(g.ev.prod(id,f))
\end{verbatim}
\end{quote}

We have defined three right objects (or functors) which makes the category
cartesian closed.  Now we define some left objects.  The object \verb"nat"
the elements of which are natural numbers can de defined as follows:
\begin{quote}
\begin{verbatim}
left object nat with pr is
  0: 1 -> nat
  s: nat -> nat
end object
\end{verbatim}
\end{quote}
The universal property in the previous section gives the usual
characterization of a natural number object in category theory.  For any
 \verb"f:1->a" and \verb"g:a->a", \verb"pr(f,g)" is from \verb"nat" to
 \verb"a", and the following equations hold:
\begin{quote}
\begin{verbatim}
pr(0,s) = id
pr(f,g).0 = f
pr(f,g).s = g.pr(f,g)
h.pr(f,g) = pr(h.f,k)
\end{verbatim}
\end{quote}
where \verb"k" is a morphism satisfying \verb"h.g = k.h".  As is well known,
\verb"pr" gives the power of defining primitive recursive functions.  For
example, a morphism which adds two natural numbers can be defined by:
\begin{quote}
\begin{verbatim}
ev.pair(pr(curry(pi2),curry(s.ev)).pi1,pi2)
\end{verbatim}
\end{quote}

Similarly, we can define any non-equational initial algebras.  For example,
the object \verb"list(a)" an element of which is a list of \verb"a" elements
can be defined by:
\begin{quote}
\begin{verbatim}
left object list(a) with prl is
  nil: 1 -> list
  cons: prod(a,list) -> list
end object
\end{verbatim}
\end{quote}
In initial algebraic specification methods, `list' is usually declared with
list destructors like `head' and `tail' and they should be connected to the
constructors `nil' and `cons' by equations.  In CPL, `head' and `tail' can
be introduced by using factorizer \verb"prl" (short for primitive recursion
on list).
\begin{quote}
\begin{verbatim}
head=prl(in2,in1.pi1)
tail=case(in1.pi2,in2).prl(in2,in1.pair(pi1,case(cons,nil).pi2))
\end{verbatim}
\end{quote}

As we said in the previous section that CPL is symmetric, we can define any
non-equational final co-algebras.  As an example, we define the object
 \verb"inf_list(a)" an element of which is an infinite list of \verb"a"
elements.
\begin{quote}
\begin{verbatim}
right object inf_list(a) with fold is
  hd: inf_list -> a
  tl: inf_list -> inf_list
end object;
\end{verbatim}
\end{quote}
For any \verb"f:b->a" and \verb"g:b->b", \verb"fold(f,g)" is from \verb"b"
to \verb"inf_list(a)" satisfying the following equations.
\begin{quote}
\begin{verbatim}
fold(hd,tl) = id
hd.fold(f,g) = f
tl.fold(f,g) = fold(f,g).g
fold(f,g).h = fold(f.h,k)    where g.h = h.k
\end{verbatim}
\end{quote}

As the dual of \verb"nat", a co-natural number object is given by:
\begin{quote}
\begin{verbatim}
right object co_nat with co_pr is
  p: co_nat -> coprod(1,co_nat)
end object
\end{verbatim}
\end{quote}
where \verb"coprod" gives coproduct of two objects defined as the dual of
 \verb"prod".  Unit morphism \verb"p" gives a predecessor function of
 \verb"co_nat".  For any \verb"f:a->coprod(1,a)", \verb"co_pr(f)" is from
 \verb"a" to \verb"co_nat" satisfying equations:
\begin{quote}
\begin{verbatim}
co_pr(p) = id
p.co_pr(f) =  co_prod(id,co_pr(f)).f
co_pr(f).g = co_pr(h)        where f.g = coprod(id,g).h
\end{verbatim}
\end{quote}


\section{Reduction Rules}

So far we have only seen the declaration mechanism for data types in CPL.
What is a program in CPL?  A CPL program is given as a morphism (or a term,
which is its syntactic representation) which is a composition of unit
morphisms, factorizers and functors (in fact functors are automatically
expanded into the equivalent terms involving factorizers by the system).
Then, what does CPL calculate?  Given a term, CPL simplifies it by using the
reduction rules which will be explained in this section.

Declaring objects in CPL gives equations which should be satisfied in a
model.

\begin{displaymath}
\begin{tabular}{ll}
\begin{tabular}{l}
right object $O(a)$ with $\psi$ is \\
\makebox[2em]{} $u$ : $F(a,O)$ $\longrightarrow$ $G(a,O)$ \\
end object \\
\end{tabular}
&
\begin{tabular}{l}
left object $\hat{O}(a)$ with $\hat{\psi}$ is \\
\makebox[2em]{} $\hat{u}$ : $\hat{F}(a,\hat{O})$ $\longrightarrow$
	$\hat{G}(a,\hat{O})$ \\
end object \\
\end{tabular}
\\
\begin{tabular}{l}
$\psi (u) = id_{O(a)}$ \\
$u\circ F(a,\psi (f)) = G(a,\psi (f))\circ f$ \\
$\psi (f)\circ g = \psi (h)$ \\
\makebox[5em]{} where $f\circ F(a,g) = G(a,g)\circ h$ \\
$O(f) = \psi (G(f,O(a))\circ u\circ F(f,O(a))$ \\
\end{tabular}
&
\begin{tabular}{l}
$\hat{\psi}(\hat{u}) = id_{\hat{O}(a)}$ \\
$\hat{G}(a,\hat{\psi}(f))\circ \hat{u} = f\circ \hat{F}(a,\hat{\psi}(f))$ \\
$g\circ \hat{\psi}(f) = \hat{\psi}(h)$ \\
\makebox[5em]{} where $\hat{G}(a,g)\circ f = h\circ \hat{F}(a,g)$ \\
$\hat{O}(f) = \hat{\psi}(\hat{G}(f,\hat{O}(a))\circ \hat{u}\circ
	\hat{F}(f,\hat{O}(a))$ \\
\end{tabular}
\end{tabular}
\end{displaymath}

From these equations, CPL uses the following reduction rules to simplify
terms.
\begin{displaymath}
\begin{array}{c}
<g,c> \Longrightarrow c' \makebox[4em]{} <f,c'> \Longrightarrow c'' \\
\hline
<f\circ g,c> \Longrightarrow c''\\
\end{array}
\end{displaymath}
\begin{displaymath}
\begin{array}{c}
<id,c> \Longrightarrow c \\
\end{array}
\end{displaymath}
\begin{displaymath}
\begin{array}{c}
<\hat{u},c> \Longrightarrow \hat{u}\circ c \\
\end{array}
\end{displaymath}
\begin{displaymath}
\begin{array}{c}
<\psi (f),c> \Longrightarrow \psi(f)\circ c \\
\end{array}
\end{displaymath}
\begin{displaymath}
\begin{array}{c}
<f\circ \hat{F}(id,\hat{\psi}(f)),c> \Longrightarrow c' \\
\hline
<\hat{G}(id,\hat{\psi}(f)),\hat{u}\circ c> \Longrightarrow
	c' \\
\end{array}
\end{displaymath}
\begin{displaymath}
\begin{array}{c}
<G(id,\psi (f))\circ f,c> \Longrightarrow c' \\
\hline
<u,F(id,\psi (f))\circ c> \Longrightarrow c' \\
\end{array}
\end{displaymath}

The reduction rules try to reduce a term to a composition of co-units (unit
morphisms associated with left objects) and right factorizers.  In order to
apply the last two reduction rules, $F$ and $\hat{G}$ should be restricted
to `productive' and `coproductive', respectively, in the second argument,
where a right object $O(a)$ is productive in $a$ if $G(a,O)$ is simply $a$
and $F(a,O)$ does not depend on $a$ and is productive in $O$.
 \verb"prod(a,b)" is productive in both arguments and \verb"exp(a,b)" is
productive in \verb"b".  The condition being `coproductive' can be defined
dually.

The reduction rules cannot reduce an arbitrary term to a canonical one, but
they do when the given term is ground (i.e.\ a term form the terminal
object) and its codomain corresponds to an algebra (e.g.\ natural number).


\section{Examples of Reduction}

Let us see some reductions done in CPL.  We have defined a morphism adding
two natural numbers.  Let us try $1+2$.
\begin{quote}
\begin{verbatim}
cpl>let add=ev.pair(pr(curry(pi2),curry(s.ev)).pi1,pi2)
add : prod(nat,nat) -> nat defined
cpl>simp add.pair(s.0,s.s.0)
0:add.pair(s.0,s.s.0)*
1:add*pair(s.0,s.s.0)
2:ev.pair(pr(curry(pi2),curry(s.ev)).pi1,pi2)*pair(s.0,s.s.0)
3:ev*pair(pr(curry(pi2),curry(s.ev)).pi1,pi2).pair(s.0,s.s.0)
4[1]:pr(curry(pi2),curry(s.ev)).pi1.pair(s.0,s.s.0)*
5[1]:pr(curry(pi2),curry(s.ev)).pi1*pair(s.0,s.s.0)
6[1]:pr(curry(pi2),curry(s.ev)).s.0*id
7[1]:pr(curry(pi2),curry(s.ev)).s*0
8[1]:pr(curry(pi2),curry(s.ev))*s.0
9[1]:curry(s.ev).pr(curry(pi2),curry(s.ev))*0
10[1]:curry(s.ev).curry(pi2).!*
11[1]:curry(s.ev).curry(pi2)*!
12[1]:curry(s.ev)*curry(pi2).!
13[1]:*curry(s.ev).curry(pi2).!
14:s.ev*pair(curry(pi2).!,pi2.pair(s.0,s.s.0))
15[1]:curry(pi2).!*
16[1]:curry(pi2)*!
17[1]:*curry(pi2).!
18:s.pi2*pair(!,pi2.pair(s.0,s.s.0))
19:s.pi2.pair(s.0,s.s.0)*id
20:s.pi2*pair(s.0,s.s.0)
21:s.s.s.0*id
22:s.s.s*0
23:s.s*s.0
24:s*s.s.0
25:*s.s.s.0
s.s.s.0
    :1 -> nat
\end{verbatim}
\end{quote}
It took 25 steps of applying the reduction rules to get the final result,
$3$.  At each step CPL prints out trace information.  The right hand side
of \verb"*" is a reduced term and the left hand side is one to be reduced.
Step 4 to 13 is a sub-reduction to get a canonical form for productive
objects.  The nesting level of reduction is printing inside brackets.

Let us see some reductions for \verb"list(a)".
\begin{quote}
\begin{verbatim}
cpl>show seq
pi2.pr(pair(0,nil),pair(s.pi1,cons))
    : nat -> list(nat)
cpl>simp seq.s.s.s.0
0:seq.s.s.s.0*
1:seq.s.s.s*0
2:seq.s.s*s.0
3:seq.s*s.s.0
4:seq*s.s.s.0
5:pi2.pr(pair(0,nil),pair(s.pi1,cons))*s.s.s.0
6:pi2.pair(s.pi1,cons).pr(pair(0,nil),pair(s.pi1,cons))*s.s.0
7:pi2.pair(s.pi1,cons).pair(s.pi1,cons).pr(pair(0,nil),pair(s.pi1,cons))*
    s.0
8:pi2.pair(s.pi1,cons).pair(s.pi1,cons).pair(s.pi1,cons).pr(pair(0,nil),
    pair(s.pi1,cons))*0
9:pi2.pair(s.pi1,cons).pair(s.pi1,cons).pair(s.pi1,cons).pair(0,nil).!*
10:pi2.pair(s.pi1,cons).pair(s.pi1,cons).pair(s.pi1,cons).pair(0,nil)*!
11:pi2.pair(s.pi1,cons).pair(s.pi1,cons).pair(s.pi1,cons)*pair(0,nil).!
12:pi2.pair(s.pi1,cons).pair(s.pi1,cons)*pair(s.pi1,cons).pair(0,nil).!
13:pi2.pair(s.pi1,cons)*pair(s.pi1,cons).pair(s.pi1,cons).pair(0,nil).!
14:pi2*pair(s.pi1,cons).pair(s.pi1,cons).pair(s.pi1,cons).pair(0,nil).!
15:cons*pair(s.pi1,cons).pair(s.pi1,cons).pair(0,nil).!
16:*cons.pair(s.pi1,cons).pair(s.pi1,cons).pair(0,nil).!
cons.pair(s.pi1,cons).pair(s.pi1,cons).pair(0,nil).!
    :1 -> list(nat)
cpl>show head
prl(in2,in1.pi1)
    : list(*a) -> coprod(*a,1)
cpl>simp head.it
0:head.it*
1:head.cons.pair(s.pi1,cons).pair(s.pi1,cons).pair(0,nil).!*
2:head.cons.pair(s.pi1,cons).pair(s.pi1,cons).pair(0,nil)*!
3:head.cons.pair(s.pi1,cons).pair(s.pi1,cons)*pair(0,nil).!
4:head.cons.pair(s.pi1,cons)*pair(s.pi1,cons).pair(0,nil).!
5:head.cons*pair(s.pi1,cons).pair(s.pi1,cons).pair(0,nil).!
6:head*cons.pair(s.pi1,cons).pair(s.pi1,cons).pair(0,nil).!
7:prl(in2,in1.pi1)*cons.pair(s.pi1,cons).pair(s.pi1,cons).pair(0,nil).!
8:in1.pi1.pair(pi1,prl(in2,in1.pi1).pi2)*pair(s.pi1,cons).pair(s.pi1,cons).
    pair(0,nil).!
9:in1.pi1*pair(pi1,prl(in2,in1.pi1).pi2).pair(s.pi1,cons).pair(s.pi1,cons).
    pair(0,nil).!
10:in1.pi1*pair(s.pi1,cons).pair(s.pi1,cons).pair(0,nil).!
11:in1.s.pi1*pair(s.pi1,cons).pair(0,nil).!
12:in1.s.s.pi1*pair(0,nil).!
13:in1.s.s.0*!
14:in1.s.s*0.!
15:in1.s*s.0.!
16:in1*s.s.0.!
17:*in1.s.s.0.!
in1.s.s.0.!
    :1 -> coprod(nat,1)
\end{verbatim}
\end{quote}

Both \verb"nat" and \verb"list(a)" are initial algebras.  Let us see how CPL
simplifies a ground term the codomain of which is a final co-algebra.
 \verb"inf_seq" gives an infinite list the elements of which are the same.
\begin{quote}
\begin{verbatim}
cpl>show inf_seq
fold(id,id)
    : *a -> inf_list(*a)
cpl>simp inf_seq.s.s.0
0:inf_seq.s.s.0*
1:inf_seq.s.s*0
2:inf_seq.s*s.0
3:inf_seq*s.s.0
4:fold(id,id)*s.s.0
5:*fold(id,id).s.s.0
fold(id,id).s.s.0
    :1 -> inf_list(nat)
cpl>let l=it
l : 1 -> inf_list(nat) defined
cpl>simp hd.l
0:hd.l*
1:hd.it*
2:hd.fold(id,id).s.s.0*
3:hd.fold(id,id).s.s*0
4:hd.fold(id,id).s*s.0
5:hd.fold(id,id)*s.s.0
6:hd*fold(id,id).s.s.0
7:id*s.s.0
8:*s.s.0
s.s.0
    :1 -> nat
cpl>simp hd.tl.l
0:hd.tl.l*
1:hd.tl.fold(id,id).s.s.0*
2:hd.tl.fold(id,id).s.s*0
3:hd.tl.fold(id,id).s*s.0
4:hd.tl.fold(id,id)*s.s.0
5:hd.tl*fold(id,id).s.s.0
6:hd.fold(id,id)*s.s.0
7:hd*fold(id,id).s.s.0
8:id*s.s.0
9:*s.s.0
s.s.0
    :1 -> nat
cpl>show inf_inc_seq
fold(id,s).0
    : 1 -> inf_list(nat)
cpl>simp hd.inf_inc_seq
0
    :1 -> nat
cpl>simp hd.tl.inf_inc_seq
s.0
    :1 -> nat
\end{verbatim}
\end{quote}

\end{document}
